\documentclass[10pt]{article}
\usepackage[a4paper,margin=1in]{geometry}
\usepackage{xeCJK}
\setCJKmainfont{FandolSong-Regular.otf}[ItalicFont=FandolKai-Regular.otf]
\setCJKsansfont{FandolHei-Regular.otf}
\setCJKmonofont{FandolFang-Regular.otf}
\usepackage{hyperref}
\usepackage{enumitem}
\setlist{nosep,leftmargin=1.5em}
\hypersetup{colorlinks=true,linkcolor=blue,urlcolor=blue,citecolor=blue}
\XeTeXlinebreaklocale "zh"
\XeTeXlinebreakskip = 0pt plus 1pt
\setlength{\parindent}{2em}
\setlength{\parskip}{0.35em}
\emergencystretch=3em
\sloppy

\title{DriveDreamer: Towards Real-world-driven World Models for Autonomous Driving\\[0.5em]\large 中文重构译文}
\author{中文译文备份}
\date{}

\begin{document}
\maketitle

\noindent\textbf{说明：} 本文档由 PDF 抽取文本重构生成，用于在缺少可用 arXiv LaTeX 源码时备份中文译文。章节和段落为自动恢复，图表、公式和双栏排版未按原始论文完整复刻。

\bigskip

\section*{DriveDreamer：面向自动驾驶的真实世界驱动世界模型}

Xiaofeng Wang*，Zheng Zhu*，Guan Huang，Xinze Chen，Jiagang Zhu，Jiwen Lu

\section*{摘要}
世界模型在自动驾驶中的用途主要体现在两方面：一是建模驾驶环境，二是生成可用于训练或评估的驾驶视频和策略信号。已有相关工作多集中在游戏或仿真环境，较少直接处理真实道路数据。本文提出 DriveDreamer，一个面向真实驾驶场景的世界模型。

真实驾驶场景包含复杂交通结构，直接在像素空间搜索代价较高。为此，论文利用扩散模型学习驾驶环境表示，并采用两阶段训练流程：第一阶段训练 Autonomous-driving Diffusion Model (Auto-DM)，把交通结构信息作为中间条件以提高采样效率，并区分动态前景目标和静态背景；第二阶段通过视频预测建立世界模型，使用驾驶动作迭代更新未来交通结构，从而根据不同驾驶策略预测未来驾驶环境。

作者在 nuScenes 基准上实例化并评估 DriveDreamer。实验结果显示，该模型能够生成可控且符合交通约束的驾驶视频，并能基于历史观测和 Auto-DM 特征预测未来驾驶策略。

\section{导言}
受通用人工智能和具身智能研究启发，自动驾驶系统正在从感知和规划模块的组合，转向更强调环境建模和未来预测的范式。世界模型适合这一问题：它们可以生成多样的驾驶视频，包括长尾场景，并为驾驶感知方法提供训练数据；同时，世界模型的预测能力也可以服务于端到端驾驶策略。

从视觉信号中学习世界模型的潜在动力学最早主要出现在视频预测研究中。通过从已观测视觉序列外推，视频预测方法可以推断环境未来状态，并建模场景中目标随时间变化的方式。然而，直接在像素空间建模复杂驾驶场景代价很高，采样空间也很大。DriveDreamer 试图通过结构化交通条件和扩散模型缓解这一问题。

\subsection{相关技术}
扩散模型通过逐步向数据加入噪声，并学习反向去噪过程来生成样本。近年来，这类模型已经被用于图像合成、视频生成和三维内容生成。为了提高可控性，ControlNet、GLIGEN、T2I-Adapter 和 Composer 等方法进一步将深度图、语义分割图、文本提示或其他结构条件注入生成过程。自动驾驶场景中的生成任务也开始利用 BEV 布局、HDMap 和 3D 边界框等交通结构信息，例如 MagicDrive、DrivingDiffusion 和 BEVControl 等工作。它们说明，结构化条件对于生成符合道路拓扑和交通约束的驾驶图像很重要。

视频生成和视频预测与世界模型关系密切。视频生成关注从条件输入生成连续帧，视频预测则根据历史观测推断未来视觉变化。已有方法包括 VAE、自回归模型、流模型、GAN 以及扩散模型。DriveGAN 等工作尝试在驾驶场景中建立可控神经仿真器，但多数相关研究仍停留在图像或视频生成层面，或者主要面向模拟环境。相比之下，DriveDreamer 试图在真实驾驶数据上建立一个能够理解交通结构、响应驾驶动作并预测未来视频的世界模型。

\subsection{总体框架}
DriveDreamer 采用两阶段训练流程。第一阶段训练 Auto-DM，使扩散模型能够在结构化交通条件下生成驾驶图像和视频。该阶段分为两个步骤：步骤 1 使用单帧结构化条件和图像监督，让模型学习 HDMap、3D 边界框与图像内容之间的对应关系；步骤 2 加入视频监督和时间注意力，使模型学习驾驶场景中的时序一致性和运动变化。

第二阶段引入 ActionFormer 建立可交互的驾驶世界模型。给定参考帧 $I_0$、道路结构信息（HDMap $H_0$ 和 3D 边界框 $B_0$）以及历史驾驶动作，ActionFormer 迭代预测未来交通结构特征，并将这些特征作为条件提供给 Auto-DM，从而生成未来驾驶视频。文本提示用于调整场景风格，例如天气和时间；历史驾驶动作与 Auto-DM 的多尺度特征则用于预测未来驾驶动作。通过这种设计，DriveDreamer 可以从静态驾驶图像生成扩展到面向不同驾驶策略的未来环境预测。

图2概括了 DriveDreamer 的两阶段训练流程。第一阶段学习交通结构约束和视频生成能力；第二阶段把驾驶动作纳入预测过程，使模型能够与环境状态演化相联系。图3展示总体框架：参考图像、HDMap、3D 边界框、文本提示和驾驶动作共同作为条件，分别影响未来交通结构、视频生成和动作预测。

\subsection{第一阶段训练}

在 DriveDreamer 中，Auto-DM 用于从真实驾驶视频中建模并理解驾驶场景。仅从像素空间理解现实驾驶场景会面对很大的搜索空间，因此作者显式引入结构化交通信息作为条件输入，以提升可控性和训练效率。

图4. Auto-DM 的整体结构。Auto-DM 接收三类控制条件：空间对齐条件（如 HDMap）会与扩散噪声特征拼接；位置条件由 3D 边界框及其类别表示，先展平后通过门控自注意力注入视觉特征；文本提示通过交叉注意力进入扩散步骤，用于控制生成驾驶视频的风格。模型还加入时间注意力层，以保持生成视频帧之间的一致性。训练时，扩散模型预测加入到潜在特征中的噪声，并通过噪声预测损失优化 Auto-DM。

为了让 Auto-DM 理解驾驶动态，作者在原有 UNet 特征上加入时间注意力。时间位置编码由正弦函数生成，视觉特征在时间维度上重排后执行自注意力，再恢复到原始维度。这样可以让模型学习帧间动态，并保持特征结构完整。同一结构也可以扩展到多视图图像和多视图视频生成：在多视图设置中，注意力主要关注相邻视角，从而加强视角一致性。

第一阶段训练包含两个步骤。第一步只使用单帧结构化交通条件作为输入，并用单帧图像监督，使模型先学习交通结构约束。HDMap 和 3D 边界框可以来自人工标注，也可以来自已有感知模型，例如 LAV、BEVerse 和 UniAD。三通道 HDMap 包含车道边界、车道分隔线和人行横道；3D 边界框通过八个角点投影到图像平面，生成对应条件。该步骤省略时间注意力层，使网络集中学习交通结构，并加快收敛。

第二步在第一步预训练模型的基础上加入时间注意力，并用驾驶视频监督。与单帧训练相比，视频训练使 Auto-DM 能够建模驾驶场景中的时序变化。第一、第二步使用与底层图像扩散模型相同的噪声调度：前向过程向潜在特征逐步加入噪声，训练目标是预测所加入的噪声。可训练参数主要位于门控自注意力、时间注意力和交叉注意力层中。

图5. ActionFormer 的整体结构。初始结构条件 $H_0$ 和 $B_0$ 先被编码并展平为一维潜在表示；随后通过自注意力和 MLP 聚合，得到隐藏状态。交叉注意力层建立隐藏状态与驾驶动作之间的关联，GRU 块则用于迭代预测未来隐藏状态。预测得到的隐藏状态会与动作特征拼接，并解码为未来交通结构条件，再作为 Auto-DM 的输入。

第二阶段训练用于建立可预测、可交互的驾驶世界模型。任务输入包括初始观测 $I_0$、$H_0$、$B_0$ 以及历史驾驶动作，目标是生成未来驾驶视频和未来驾驶动作。已训练的 Auto-DM 可以根据结构化交通信息序列生成驾驶视频，但视频预测时未来交通结构条件不可直接获得。ActionFormer 解决这一问题：它利用历史驾驶动作迭代预测未来结构条件，并在特征层面预测未来交通结构，从而减少像素级噪声干扰，使预测更稳健。

在概率建模上，DriveDreamer 将未来视频预测和未来动作预测联合起来。由于隐藏状态更新由 GRU 确定，推断时主要需要估计潜在变量。作者引入变分分布进行近似推断，并使用简化的变分下界训练模型。实践中，视频预测部分使用均方误差损失，动作预测部分使用 L1 损失。相关可学习层包括 ActionFormer、Auto-DM、视频解码器（VAE 解码器）和动作解码器。

训练实现上，Auto-DM 基于 Stable Diffusion v1.4 构建，并冻结其原始参数。第一阶段的第一步训练 40 个 epoch，批量大小为 16；第二步训练 10 个 epoch，批量大小为 1，视频长度 $N=32$，空间分辨率为 $448\times256$。第二阶段视频预测训练中，模型预测 16 帧驾驶视频以及 16 个未来驾驶动作，同样训练 10 个 epoch，批量大小为 1。所有实验在 A800 GPU 上完成，优化器为 AdamW，学习率为 $5\times10^{-5}$。

评价方面，作者同时进行定性与定量评估。生成质量使用逐帧 Fréchet Inception Distance（FID）和 Fréchet Video Distance（FVD）衡量，评估图像统一缩放到 $448\times256$。为验证生成图像对驾驶感知训练的帮助，实验还在 3D 目标检测任务上评价 DriveDreamer，基线为 FCOS3D 和 BEVFusion。此外，作者按照 ST-P3 的设置评价输出驾驶轨迹，以检验驾驶动作生成能力。

在未来驾驶动作预测中，DriveDreamer 首先从 Auto-DM 中汇聚多尺度 UNet 特征，并将这些视觉特征与历史动作特征拼接，再通过 MLP 层解码得到未来 3 秒的驾驶动作。

\subsection{可控驾驶视频生成}

经过两阶段训练后，DriveDreamer 可以利用交通结构约束预测未来驾驶状态，并把动作输入纳入视频生成过程。本节展示 DriveDreamer 在结构化交通条件下生成多样驾驶视频的能力，并验证生成图像能否增强驾驶感知方法的训练。同时，模型还能响应不同输入动作，从而控制车辆轨迹并生成不同驾驶视频。

如图1和图6所示，DriveDreamer 可以生成严格符合结构化交通条件的图像和视频。生成条件包括 HDMap 和 3D 边界框，来自 nuScenes 验证集；图中用红色矩形和黄色圆圈标出的区域表明，生成图像具有多视图一致性，并与真实条件对齐。除结构化交通条件外，文本提示也可以诱导天气和时间等变化，例如晴天、雨天和夜晚。

为进一步验证生成质量，作者从 nuScenes 训练集中提取 4000 组交通条件生成驾驶图像，并将生成图像与真实图像结合，用于训练 3D 目标检测任务。表1显示，加入合成数据后，FCOS3D 和 BEVFusion 的 mAP 分别提升 0.7 和 3.0，说明生成数据能够改善下游感知任务。表2比较了 nuScenes 验证集上的生成质量。即使不使用第一阶段训练，DriveDreamer 的 FID 和 FVD 也优于 DriveGAN；经过第一阶段训练后，模型对驾驶场景中的结构信息理解更好，视频生成质量进一步提升。最后，ActionFormer 能有效利用第一阶段学到的交通结构知识，并根据输入动作迭代更新未来结构信息，相比直接拼接动作特征的基线进一步改善视频质量。

除了使用结构化交通条件生成驾驶视频，DriveDreamer 还能根据不同驾驶动作产生多样视频。从同一初始帧及其结构信息出发，输入不同驾驶动作后，模型可以生成左转、右转等不同未来轨迹。这说明 DriveDreamer 可以作为一种可控数据生成工具，用于补充自主驾驶系统中的角落案例和长尾场景训练数据。

\subsection{驾驶动作生成}

除可控驾驶视频生成外，DriveDreamer 也用于预测未来驾驶动作。给定初始帧和历史驾驶动作，模型可以生成符合现实场景的未来动作。作者进一步进行定量评估：使用 MLP 层编码历史驾驶动作信息，并汇聚多尺度 UNet 特征作为视觉提示，随后拼接两种模态特征来学习未来驾驶轨迹。nuScenes 上的开环评价结果见表3。DriveDreamer 的平均 L2 轨迹误差为 0.29 m，优于多模态方法 VAD；同时，相比文献[75]报告的结果，平均碰撞率相对降低 21\%。这说明 DriveDreamer 学得的视觉特征有助于端到端自主驾驶，并能提升驾驶安全性。

\section{讨论与结论}

DriveDreamer 是面向真实世界驾驶场景的世界模型探索。它以真实驾驶数据为基础，结合扩散模型生成能力，学习交通结构、驾驶视频和动作之间的关系。以往世界模型研究更多集中在游戏或模拟环境中，而 DriveDreamer 将世界建模扩展到真实交通条件下的复杂驾驶场景。

这项工作给出了自动驾驶世界模型的一种实现路线。DriveDreamer 表明，结合结构化交通条件、文本提示和驾驶动作，可以构建兼具可控生成和预测能力的驾驶世界模型；后续工作仍需要进一步验证这类模型在闭环驾驶和安全评估中的作用。

\section*{参考文献说明}

原论文参考文献主要覆盖扩散模型、视频生成、自动驾驶感知与规划、世界模型和 nuScenes 相关基准。由于当前译文由 PDF 文本抽取重构，原始参考文献列表存在双栏错位和断行粘连问题；为避免将不可读的条目误作为正式译文，本文档不逐条复刻参考文献。需要精确引用时，建议以原论文 PDF 或官方 BibTeX 为准。


\section*{补充材料摘要}

补充材料进一步给出了 DriveDreamer 的实现细节，包括条件编码器结构、多视图生成、动作预测结构和合成数据训练流程。编码器部分分别处理参考图像、HDMap、3D 边界框和动作等条件输入；空间对齐条件通过卷积层下采样，非结构化条件则通过 MLP 编码。

DriveDreamer 可以自然扩展到多视图图像和多视图视频生成。多视图图像生成与视频生成框架基本一致，只是将帧间注意力替换为视角间注意力；多视图视频生成则同时堆叠视角间注意力和帧间注意力，从而获得视角一致、帧间一致的视频。动作预测部分会融合多尺度 UNet 特征和编码后的驾驶动作特征，并通过 MLP 学习未来驾驶动作。

合成数据训练中，作者使用 DriveDreamer 生成的数据增强 3D 检测器训练。模型先生成高分辨率图像，再调整到原始分辨率用于 FCOS3D 和 BEVFusion 等检测器。可视化结果表明，模型既能根据结构化交通条件生成多样驾驶场景，也能通过文本提示控制天气和时间，并在交叉路口、交通灯和转弯等复杂情形下预测合理的驾驶动作。

\section*{补充可视化说明}

补充图 9 展示了在结构化交通条件（HDMap 和 3D 边界框）下的驾驶视频生成结果。文本提示用于调节场景风格，例如晴天、雨天或夜晚。补充图 10 展示了动作条件的作用：在相同初始条件下，不同驾驶动作（如左转、右转）会生成对应的未来驾驶视频。补充图 11 将模型预测的未来驾驶动作与真实驾驶视频进行对照，用于说明动作预测与视觉结果的一致性。


\end{document}
